\clearpage
\section{Inclusion and Exclusion Process}
\label{app:schema}

We conducted a systematic review, using a combination of keyword search, LLM filtering, and human filtering to identify articles. Figure~\ref{fig:schema} shows the steps of the search process and the number of papers included at each step.

\begin{figure}[t!]
    \centering
    \includegraphics[width=\linewidth]{images/inclusion_exclusion.png}
    \caption{\textbf{Flowchart of the systematic review process.} Searching across EMNLP, NAACL, ACL, ICML, ICLR, and NeurIPS, we identified 2,189 papers matching the keyword search, and 445 which ultimately met the review criteria.}
    \label{fig:schema}
\end{figure}

\subsection{Keyword Search and LLM Filtering}

Keyword search across the six target conferences identified 2,189 papers which included the keywords `benchmark' and `LLM' or `language model' in the title or abstract. A manual scan of these articles indicated that many articles were technical papers developing techniques for language modelling which reported improvements on various benchmarks. Since these papers were not the target of the review, further filtering was conducted via LLM.

For the LLM filtering, we used GPT-4o mini to identify the features about the papers relevant to inclusion and exclusion. We processed the inclusion criteria in order, removing articles at each step which did not meet the criteria. Table~\ref{tab:llm_exclusion} shows the steps of this process.

\begin{table}[h!]
    \centering
    \renewcommand{\arraystretch}{1.3}
    {\small
    \begin{tabular}{p{.2\textwidth}p{.45\textwidth}p{.1\textwidth}p{.1\textwidth}}
            \toprule
         \textbf{\newline Criterion} & \textbf{\newline LLM Prompt} & \textbf{Articles\newline Excluded} & \textbf{Articles\newline Remaining} \\
         \midrule
         System Prompt & You are an academic assistant, filtering articles to identify which ones are relevant to a literature review. & - & 2,189\\
         Exclude articles which do not implement a benchmark & Please read the paper title and abstract below, and tell me whether the paper creates and describes a new benchmark for large language models.
                  
         After reading the title and abstract, please very briefly describe whether the article implements a new benchmark for large language models. Then, on a new line write **Answer:** followed by a single word answer of `Yes' or `No' as to whether the article creates and describes a new benchmark. & 1,251 & 938 \\
         Exclude articles which require modalities beyond text and vision & Please read the paper title and abstract below, and tell me the primary modality of the dataset being used.
                  
        After reading the title and abstract, please very briefly describe the primary modality of the article. Then, on a new line write **Answer:** followed by a single word answer describing the primary modality considered in the article. Your answer should be either Language, Image, Video, Audio, Multimodal or Other. Use Other only when the primary modality is not one of the previous options. & 92 & 846 \\
        Exclude articles which are not primarily about creating a new benchmark & Please read the paper title and abstract below, and tell me the primary focus area of the paper.
                  
        After reading the title and abstract, please very briefly describe the primary focus of the paper. Then, on a new line write **Answer:** followed by a single word answer of `Benchmark', `Technical', `Methodological' or `Other' to categorize the primary contribution. & 324 & 522 \\
        \bottomrule
        \\
    \end{tabular}
    \caption{\textbf{LLM Filtering Steps.} The prompts used for progressive filtering of the articles to be reviewed, and the number of articles excluded at each step.}
    \label{tab:llm_exclusion}
    }
\end{table}

We validated the results of this process by randomly selecting 50 articles from the 2,189 and manually categorising them for inclusion and exclusion. Table~\ref{tab:exclusion_validation} shows the confusion matrix of the LLM exclusion relative to the human gold standard classes. The overall precision in this subset is 80\% and the recall is 89\%. As such we expect the filtering to be highly effective in capturing the relevant articles for human review, though not perfect, which we note as a limitation. 

Of the overall batch of 2,189, 522 articles were selected for review, about 24\%, which is similar to the 20\% of articles selected in the random sample. The fraction of the manually reviewed articles which were included in the final study was 445, indicating that the precision in the full study was about 85\%, similar to this sample. 

\begin{table}[h!]
    \centering
    \renewcommand{\arraystretch}{1.3}
{\small
    \begin{tabular}
    {p{.3\textwidth}p{.1\textwidth}p{.1\textwidth}p{.1\textwidth}p{.1\textwidth}p{.1\textwidth}}
    \toprule
        \textbf{\newline Criterion} & \textbf{Articles \newline Compared} & \textbf{True\newline Inclusion} & \textbf{False\newline Inclusion} & \textbf{True\newline Exclusion} & \textbf{False\newline Exclusion}\\
        \midrule
        Exclude articles which do not implement a benchmark & 50 & 14 & 9 & 22 & 5\\
        Exclude articles which require modalities beyond text and vision & 14 & 13 & 0 & 1 & 0\\
        Exclude articles which are not primarily about creating a new benchmark. & 9 & 8 & 0 & 0 & 1\\
        Overall & 50 & 8 & 2 & 39 & 1 \\ 
        \bottomrule
        \\
    \end{tabular}
    \caption{\textbf{LLM Filtering Human Validation.} A comparison between human and LLM filtering on a subset of 50 articles drawn at random from the initial list of 2,189. Human filtering results are treated as gold-standard. The steps were conducted in the same order as the real filtering, so that the number of papers remaining falls at each step. The confusion matrix is shown for each step of the filtering process.}
    \label{tab:exclusion_validation}
    }
\end{table}



